\documentclass{article}
\usepackage{ctex}
\usepackage{amsmath}
\usepackage{tikz}
\usepackage{unicode-math}
\usetikzlibrary{arrows.meta, decorations.markings,patterns}
\begin{document}

2.1-5 %三平行金属板$ A, B$和 $C$,面积都是 $200 cm^2,A,B$ 板相距$4.0$ mm,$A,C$ 板相距 $2.0$ mm,$B,C$ 两板都接地.
%如果使$A$板带正电$ 3.0\times10^{-7}$ C,在略去边缘效应时,问 $B$ 板和 $C$ 板上感应电荷各是多少?以地的电势为零,问 $A $板的电势是多少?


$$ S = 200\ \rm{cm}^2 = 200 \times 10^{-4} \rm{m}^2 = 0.02\rm{m}^2 $$

$ A $ 带正电荷 $ Q_A = 3.0 \times 10^{-7} \rm{C} $

$ A $ 与 $ B $ 之间、$ A $ 与 $ C $ 之间各形成一个电容器

因为 $ B $ 和 $ C $ 接地,$U_B=U_C=0$

$$U_A = U_{AB} = U_{AC}$$

平行板电容公式：
$$C = \frac{\varepsilon_0 S}{d},\varepsilon_0 = 8.85 \times 10^{-12} \ \rm{F/m}$$
$$C_{AB} = \frac{\varepsilon_0 S}{d_{AB}} = \frac{8.85\times 10^{-12} \times 0.02}{0.004}= 4.425\times 10^{-11} \ \rm{F}$$
$$C_{AC} = \frac{\varepsilon_0 S}{d_{AC}} = \frac{8.85\times 10^{-12} \times 0.02}{0.002}= 8.85\times 10^{-11} \ \rm{F}$$


$ A $ 板总电荷 $ Q_A $ 分布在两个表面上：  

面对 $ B $ 的表面电荷为 $ Q_{AB} $,面对 $ C $ 的表面电荷为 $ Q_{AC} $,且：
$$Q_{AB} + Q_{AC} = Q_A ,Q_{AB} = C_{AB} U_A, Q_{AC} = C_{AC} U_A$$

$$(C_{AB} + C_{AC}) U_A = Q_A$$
$$U_A = \frac{Q_A}{C_{AB} + C_{AC}}$$


$$C_{AB} + C_{AC} = 1.3275\times 10^{-10} \ \rm{F},U_A = \frac{3.0\times 10^{-7}}{1.3275\times 10^{-10}} \approx 2.26\times 10^{3} \rm{V}$$

$ B $ 板上的感应电荷 $ Q_B $
$$Q_B = - Q_{AB} = - C_{AB} U_A= - C_{AB}\times \frac{Q_A}{C_{AB} + C_{AC}}  =- 1.0\times 10^{-7} \ \rm{C}$$

同理,$ C $ 板感应电荷 $ Q_C $：
$$Q_C = - Q_{AC} = - C_{AC} U_A = -2.0\times 10^{-7} \ \rm{C}$$

$$\boxed{Q_B = -1.0\times 10^{-7}\ \rm{C},\quad Q_C = -2.0\times 10^{-7}\ \rm{C},\quad U_A \approx 2.26\times 10^{3}\ \rm{V}}$$

















\newpage

2.1-6 %点电荷 $q$ 处在导体球壳的中心,壳的内外半径分别为 $R_1$ 和 $R_2$. 求场强和电势的分布,并画出 $E-r$ 和 $U-r$ 曲线.

球壳内半径 $ R_1 $，外半径 $ R_2 $。
由于球对称，所有物理量只与 $ r $ 有关。


导体内部（ $ R_1 < r < R_2 $）电场为零

因为导体内部场为零，高斯面在 $ R_1 < r < R_2 $ 内必须电通量为零

所以内表面总电荷 = $-q$

中心点电荷 $ q $ 在内表面会感应出均匀分布的电荷 $ -q $。  

球壳原本总电荷为 0，所以外表面均匀分布总电荷 = $ +q $

$$\rho_{\text{内}}(r=R_1) = -\frac{q}{4\pi R_1^2}, \quad\rho_{\text{外}}(r=R_2) = +\frac{q}{4\pi R_2^2}.$$



用高斯定理，取同心球面为高斯面。

(a) $ 0 < r < R_1 $
高斯面内只有中心点电荷 $ q $：
$$E \cdot 4\pi r^2 = \frac{q}{\varepsilon_0}$$
$$E(r) = \frac{q}{4\pi\varepsilon_0 r^2}$$


(b) $ R_1 < r < R_2 $
导体内部，电场为零：
$$E(r) = 0.$$

(c) $ r > R_2 $
高斯面内的总电荷 =  $ q $。
$$E(r) = \frac{q}{4\pi\varepsilon_0 r^2}.$$



$$
E(r) =
\begin{cases}
\dfrac{q}{4\pi\varepsilon_0 r^2}, & 0 < r < R_1, \\
0, & R_1 < r < R_2, \\
\dfrac{q}{4\pi\varepsilon_0 r^2}, & r > R_2.
\end{cases}
$$



取无穷远为电势零点：$ U(\infty) = 0 $。

(a) $ r \ge R_2 $
$$
U(r) = \int_{r}^{\infty} E(r') \, dr'
= \int_{r}^{\infty} \frac{q}{4\pi\varepsilon_0 r'^2} dr'
= \frac{q}{4\pi\varepsilon_0 r}.
$$

(b) $ R_1 \le r \le R_2 $
导体等势，所以：
$$
U(r) = U(R_2) = \frac{q}{4\pi\varepsilon_0 R_2}.
$$

(c) $ 0 < r < R_1 $
从 $ r $ 到 $ R_1 $ 有电场 $ E = \frac{q}{4\pi\varepsilon_0 r'^2} $

从 $ R_1 $ 到 $ R_2 $ 电场为零（电势不变），所以：

$$
U(r) = \int_{r}^{R_1} \frac{q}{4\pi\varepsilon_0 r'^2} dr' + U(R_1)= \frac{q}{4\pi\varepsilon_0} \left( \frac{1}{r} - \frac{1}{R_1} \right) + \frac{q}{4\pi\varepsilon_0 R_2}.
$$

$$
U(r) = \frac{q}{4\pi\varepsilon_0} \left( \frac{1}{r} - \frac{1}{R_1} + \frac{1}{R_2} \right).
$$



$$
U(r) =
\begin{cases}
\dfrac{q}{4\pi\varepsilon_0} \left( \dfrac{1}{r} - \dfrac{1}{R_1} + \dfrac{1}{R_2} \right), & 0 < r < R_1, \\
\dfrac{q}{4\pi\varepsilon_0 R_2}, & R_1 \le r \le R_2, \\
\dfrac{q}{4\pi\varepsilon_0 r}, & r \ge R_2.
\end{cases}
$$





\vspace{2em}

\begin{tikzpicture}[scale=1.6]
    % 坐标轴
    \draw[->] (-0.5,0) -- (6,0) node[right] {$r$};
    \draw[->] (0,-0.5) -- (0,3) node[above] {$E(r), U(r)$};

    
    % E(r) - 蓝色
    \draw[domain=0.5:1.99, smooth, samples=100, thick, blue] 
        plot (\x, {1/(\x*\x)});
    \draw[domain=2.01:3.99, smooth, samples=2, thick, blue] 
        plot (\x, 0);
    \draw[domain=4.01:5.5, smooth, samples=100, thick, blue] 
        plot (\x, {1/(\x*\x)});
    
    % U(r) - 红色
    \draw[domain=0.5:1.99, smooth, samples=100, thick, red] 
        plot (\x, {1*(1/\x - 1/4)});
    \draw[domain=2:4, smooth, samples=2, thick, red] 
        plot (\x, 1);
    \draw[domain=4:5.5, smooth, samples=100, thick, red] 
        plot (\x, {1/\x});
    
    % 参考线标注
    \draw[dashed, gray] (2,0) -- (2,3) node[below right] at (2,0){$R_1$};
    \draw[dashed, gray] (4,0) -- (4,3) node[below right] at (4,0){$R_2$};

    
    % 图例
    \draw[blue, thick] (4,2.8) -- (4.5,2.8) node[right, black] {$E(r)$};
    \draw[red, thick] (5,2.8) -- (5.6,2.8) node[right, black] {$U(r)$};
    



\end{tikzpicture}





















\newpage
2.1-7 在上题中,若 $q = 4 \times 10^{-10} \, \text{C}, R_1 = 2 \rm{cm}, R_2 = 3 \rm{cm}$, 求:  

(1) 导体球壳的电势；  

(2) 离球心 $r = 1 \rm{cm}$ 处的电势；  

(3) 把点电荷移开球心 $1 \rm{cm}$,求导体球壳的电势.


$$q = 4\times 10^{-10} \ \text{C}, \quad R_1 = 2\ \text{cm} = 0.02\ \text{m}, \quad R_2 = 3\ \text{cm} = 0.03\ \text{m}$$

$$\frac{1}{4\pi\varepsilon_0} = 9\times 10^9 \ \text{N·m}^2/\text{C}^2$$



 (1)
上一题已得：对于 $ R_1 \le r \le R_2 $，导体是等势体，电势为
$$U_1 = \frac{q}{4\pi\varepsilon_0 R_2}$$
代入：
$$U_1 = (9\times 10^9) \times \frac{4\times 10^{-10}}{0.03}= (9\times 10^9) \times (1.333\ldots\times 10^{-8})= 120 \ \text{V}$$



 (2) 

$$U(r) = \frac{q}{4\pi\varepsilon_0} \left( \frac{1}{r} - \frac{1}{R_1} + \frac{1}{R_2} \right), \quad 0 < r < R_1$$
$$r = 0.01\ \text{m}, \quad R_1 = 0.02, \quad R_2 = 0.03$$
$$\frac{1}{r} = 100, \quad \frac{1}{R_1} = 50, \quad \frac{1}{R_2} \approx 33.333\ldots$$
括号内：
$$100 - 50 + 33.333\ldots = 83.333\ldots$$
$$U = (9\times 10^9) \times (4\times 10^{-10}) \times 83.333\ldots$$
 $ 9\times 10^9 \times 4\times 10^{-10} = 3.6 $
$$U = 3.6 \times 83.333\ldots \approx 300 \ \text{V}$$




 (3) 点电荷移开球心 $ 1\ \text{cm} $，求导体球壳的电势


由于静电屏蔽，壳外的电场完全由外表面的 $+q$ 均匀分布决定

壳外电场 = 位于球心的 $+q$ 产生的电场。  

因此壳的电势（$ r \ge R_2 $）为：
$$U_1 = \frac{q}{4\pi\varepsilon_0 R_2}$$


因此：
$$U_1 = 120\ \text{V} \quad \text{同上(1)}$$





$$
(1)\ 120\ \text{V}, \quad (2)\ 300\ \text{V}, \quad (3)\ 120\ \text{V}
$$






\newpage

2.1-8半径为$R_1$的导体球带有电荷 $q$,球外有一个内,外半径为$R_2,R_3$ 的同心导体球壳,壳上带有电荷 $Q$

(1)求两球的电势$U_{1}$ 和$U_2$

(2)两球的电势差$\Delta U;$

(3)以导线把球和壳连接在一起后,$U_1,U_2$和$\Delta U$分别是多少?

(4)在情形(1),(2)中,若外球接地,$U_1,U_2$和$\Delta U$为多少?

(5) 设外球离地面很远,若内球接地,情况如何?






\newpage

2.1-12 同轴传输线是由两个很长且彼此绝缘的同轴金属直圆柱体构成.
设内圆柱体的电势为$U_{1}$,半径为 $R_1$,外圆柱体的电势为 $U_2$,内半径为 $R_2$,求其间离轴为 $r$ 处$(R_1<r<R_2)$ 的电势.











\newpage

在静电场的作用下,金属导体内部的场强在任何条件下都为零?在很薄的金属薄膜(厚度为几十个甚至几个原子层)内呢?
课本阐述到达静电平衡状态后导体内的场强$E$可以等于零的隐含意义:自由电子数量足够多！查阅材料,根据金属自身的特性(自由电子体密度)来讨论这个问题.




\end{document}